\documentclass[11pt]{article}
\usepackage[margin=1in]{geometry}
\usepackage{amsmath,amssymb,amsthm,booktabs,lmodern,microtype}
\usepackage[colorlinks=true,urlcolor=blue,citecolor=blue,linkcolor=blue]{hyperref}
\newtheorem{theorem}{Theorem}
\newcommand{\Z}{\mathbb Z}
\newcommand{\R}{\mathbb R}
\newcommand{\wt}{\widehat\tau}
\title{An explicit SU(5) lens-space counterexample\ to the printed Guadagnini--Pilo conjecture}
\author{Alec Kriebel\\Independent researcher\\\small\href{https://orcid.org/0009-0001-9320-500X}{ORCID: 0009-0001-9320-500X}}
\date{October 5, 2026}
\begin{document}
\maketitle
\begin{abstract}
For the ordinary full SU(5) Reshetikhin--Turaev invariant at WZW level
$k=5$, or shifted level $r=10$, the lens spaces $L(5,1)$ and $L(5,2)$
have fundamental group $\Z/5$ but unequal nonzero magnitudes. With
$\wt(S^3)=1$, their squared magnitudes are $3475+1550\sqrt5$ and
$4025+1800\sqrt5$. A short exact root-lattice reduction proves the
calculation, conditional on the established modular-category and surgery
formulas; finite cyclotomic programs provide reproducible checks.
This contradicts Conjecture 7.5 as printed in Ohtsuki's collection.
Historical priority is unresolved: Takahito Kuriya's directly relevant
preprint could not be obtained, and no firstness or new historical
resolution is claimed. This is an unrefereed research note with extensive
AI assistance and no conventional human peer review.
\end{abstract}

\section{Claim and conventions}
Ohtsuki's edited collection \cite{O} prints Conjecture 7.5 on p.~474,
attributed to Guadagnini and Pilo \cite{GP}: the absolute value of a
nonzero quantum $G$ invariant of a closed oriented 3-manifold is
determined by its fundamental group. The volume is nominally dated 2002
and records publication on June 1, 2004. We refer to this printed claim,
without asserting its present literature status. Its scope includes
compact simple simply connected $G=\mathrm{SU}(5)$; it is not restricted
to the SU(2) lens-space theorem in \cite{GP}.

\begin{theorem}
In the full ordinary SU(5) theory at WZW level $k=5$, shifted level
$r=k+h^\vee=10$, let $\wt(M)=\tau(M)/\tau(S^3)$. Then
\begin{align}
 |\wt(L(5,1))|^2&=3475+1550\sqrt5,\label{eq:one}\\
 |\wt(L(5,2))|^2&=4025+1800\sqrt5.\label{eq:two}
\end{align}
Both manifolds have fundamental group $\Z/5$. Both values are positive,
and their difference is $550+250\sqrt5>0$. Thus the printed nonzero-magnitude
claim fails.
\end{theorem}

We use the conventions of Hansen--Takata \cite{HT}, specifically the
available preprint arXiv:math/0209403v2. Its final journal version has not
been compared here. The invariant $\tau$ there has $\tau(S^3)=S_{00}>0$.
Their $L(p,q)$ is unknot surgery with coefficient $-p/q$. A convention
using $+p/q$ reverses orientation and conjugates the ordinary unitary
invariant, leaving the magnitudes unchanged. The quantum-group root is
$e^{\pi i/10}$; the renormalized coupling in Guadagnini--Pilo notation
corresponds to shifted $r=10$, not WZW $k=5$.

The full category has all dominant labels $(a_1,a_2,a_3,a_4)$ with
$a_i\geq0$ and $\sum_i a_i\leq5$, hence $\binom94=126$ labels. No
root-lattice restriction on those labels, projective quotient, bundle
sector, spin choice, inserted link, or retained manifold framing is used.
The root lattice below is a computational quotient in a formula for
this full category. The existence and identification of the category,
and the general RT surgery theorem, are standard imported mathematical
inputs from \cite{HT}; we do not rebuild those foundations.

\section{Exact lattice reduction}
Realize the $A_4$ root system in $H=\{x\in\R^5:\sum_i x_i=0\}$ with
root lattice $Y=H\cap\Z^5$, weight lattice $X=Y^*$, and
$\rho=(2,1,0,-1,-2)$. Roots have squared length 2,
$|\rho|^2=10$, rank is 4, and there are 10 positive roots.
The simple-root Gram determinant is 5, so $\operatorname{vol}(Y)=\sqrt5$
and $\operatorname{vol}(X)=1/\sqrt5$.

Write $\mathcal S(q/p)=12s(q,p)$ for the Dedekind symbol. For $p>0$,
$s(q,p)=\sum_{n=1}^{p-1}((n/p))((qn/p))$, with
$((x))=x-\lfloor x\rfloor-1/2$ off integers and zero at integers.
Hansen--Takata Theorem 5.1 specializes to
\begin{align}
\tau(L(p,q))={}&\frac{i^{10}e^{\pi i\mathcal S(q/p)}}{(10p)^2\sqrt5}
\sum_{w\in S_5}\det(w)e^{-2\pi i\langle\rho,w\rho\rangle/(10p)}
\sum_{\nu\in Y/pY}e^{\pi i(10q/p)|\nu|^2}
 e^{(2\pi i/p)\langle\nu,q\rho-w\rho\rangle}.
\label{eq:HT}
\end{align}
Here $\gcd(p,q)=1$ and $p\ne0$. This theorem has no coprimality
requirement between $p$ and $r$; the coprime simplification in their
Proposition 5.2 is inapplicable when $p=5,r=10$ and is not used.

Set $p=5$. Representatives of $Y/5Y$ are
\[
\nu(a)=(a_1,a_2-a_1,a_3-a_2,a_4-a_3,-a_4),\quad 0\leq a_i<5.
\]
The simple roots are a free integral basis, so this gives all $5^4=625$
classes. The quadratic phase in \eqref{eq:HT} is
$e^{2\pi iq|\nu|^2}=1$. For an integral sum-zero vector $v$, character
orthogonality gives
\begin{equation}
\sum_{\nu\in Y/5Y}e^{2\pi i\langle\nu,v\rangle/5}
=\prod_{j=1}^4\sum_{a=0}^4 e^{2\pi ia(v_j-v_{j+1})/5}
=\begin{cases}625,&v_1\equiv\cdots\equiv v_5\pmod5,\\0,&\text{otherwise.}\end{cases}
\label{eq:char}
\end{equation}
The criterion is $v\in5X$, not $v\in5Y$. Since $\rho$ exhausts the
five residues, precisely five permutations survive for each nonzero
$q\pmod5$.

Put $z=e^{2\pi i/50}$, $t=z^5=e^{\pi i/5}$, and
\[
B_q=\sum_{w:\ q\rho-w\rho\in5X}\det(w)z^{-\langle\rho,w\rho\rangle}.
\]
Table~\ref{tab:survivors} lists every survivor for $q=1,2$. Signs are
permutation parities; dot products are elementary integer calculations.
\begin{table}[ht]
\centering
\begin{tabular}{c r r r}
\toprule
$q$ & $w\rho$ & $\det(w)$ & $\langle\rho,w\rho\rangle$\\
\midrule
1 & $(2,1,0,-1,-2)$ & $+1$ & 10\\
1 & $(1,0,-1,-2,2)$ & $+1$ & 0\\
1 & $(0,-1,-2,2,1)$ & $+1$ & $-5$\\
1 & $(-1,-2,2,1,0)$ & $+1$ & $-5$\\
1 & $(-2,2,1,0,-1)$ & $+1$ & 0\\
\midrule
2 & $(2,0,-2,1,-1)$ & $-1$ & 5\\
2 & $(1,-1,2,0,-2)$ & $-1$ & 5\\
2 & $(0,-2,1,-1,2)$ & $-1$ & $-5$\\
2 & $(-1,2,0,-2,1)$ & $-1$ & 0\\
2 & $(-2,1,-1,2,0)$ & $-1$ & $-5$\\
\bottomrule
\end{tabular}
\caption{The five Weyl survivors for each target lens space.}\label{tab:survivors}
\end{table}
Thus
\begin{equation}
B_1=t^{-2}+2+2t,\qquad B_2=-(1+2t^{-1}+2t).
\label{eq:B}
\end{equation}
In the positive real embedding $t+t^{-1}=(1+\sqrt5)/2$. Consequently
\begin{equation}
|B_1|^2=11+2\sqrt5,\qquad |B_2|^2=9+4\sqrt5.
\label{eq:norms}
\end{equation}
For the first identity, expand the norm as
$9+4(t+t^{-1})+2(t^2+t^{-2})+2(t^3+t^{-3})$ and use
$t^2+t^{-2}=(\sqrt5-1)/2$ and $t^3+t^{-3}=-(t^2+t^{-2})$.
For the second, $B_2=-(2+\sqrt5)$.

\section{Normalization and topology}
At $p=1$ the quotient has one class and the Dedekind symbol is zero.
Let
\[
A=\sum_{w\in S_5}\det(w)e^{-2\pi i\langle\rho,w\rho\rangle/10}.
\]
The Weyl denominator identity gives $A=-P$, where
\[
P=\prod_{j=1}^4(2\sin(\pi j/10))^{5-j}=5\sqrt5-10>0.
\]
Indeed $(2\sin18^\circ)(2\sin54^\circ)=1$,
$(2\sin36^\circ)(2\sin72^\circ)=\sqrt5$, and
$(2\sin36^\circ)^2=(5-\sqrt5)/2$. Therefore
\begin{equation}
 A=10-5\sqrt5,\quad |A|^2=225-100\sqrt5,\quad
 s=S_{00}=\tau(S^3)=\frac{-A}{100\sqrt5}>0.
\label{eq:A}
\end{equation}
Canceling the prefactors of \eqref{eq:HT} at $p=5$ and $p=1$ gives
\begin{equation}
 \wt(L(5,q))=25e^{\pi i\mathcal S(q/5)}B_q/A,
 \qquad |\wt(L(5,q))|^2=\frac{625|B_q|^2}{225-100\sqrt5}.
\label{eq:ratio}
\end{equation}
Since $(225-100\sqrt5)(225+100\sqrt5)=625$, equations
\eqref{eq:norms} and \eqref{eq:ratio} prove \eqref{eq:one}--\eqref{eq:two}.
The denominator is positive, for $225^2>100^2\cdot5$.

For comparison with vacuum matrix words, the positive rank is
$D=s^{-1}=100+40\sqrt5$, and $s^2=(9-4\sqrt5)/2000$.
The ordinary normalized modular matrix is
\[
 S_{\lambda\mu}=-\frac1{100\sqrt5}
 \sum_{w\in S_5}\det(w)e^{-2\pi i\langle w(\lambda+\rho),\mu+\rho\rangle/10},
\quad
\theta_\lambda=e^{\pi i(|\lambda+\rho|^2-|\rho|^2)/10}.
\]
The categorical Hopf matrix is $DS$, not $S$. The two positive surgery
chains are $5=[5]$ and $5/2=[3,2]$, giving words
$S\theta^5S$ and $S\theta^3S\theta^2S$. For either word $F$, the
$S^3$-normalized magnitude is $|F_{00}|/s$, with exactly one division
by $s$. Hansen--Takata Corollary 4.2 and Eqs.~(26)--(30), (35) supply
this full-category normalization. Here central charge is
$5\cdot24/10=12$, so the linear $T$ is $-\theta$ and the anomaly
$\Delta/D=-1$. These and the signature corrections have modulus one.

Replacing $q$ by $-q$ conjugates \eqref{eq:HT}: use the longest Weyl
element, which sends $\rho$ to $-\rho$ and has sign $+1$, and
$s(-q,5)=-s(q,5)$. Thus $q=4$ gives the same magnitude as $q=1$,
and $q=3$ the same as $q=2$. Inverse-denominator conventions also
exchange 2 and 3. Neither convention removes the discrepancy. A fixed
nonzero theory-wide rescaling cannot make unequal magnitudes equal.

Finally, the unknot exterior has fundamental group generated by its
meridian $\mu$, and its preferred longitude is nullhomotopic.
Filling at $-5/q$ kills $\mu^5$, so van Kampen gives
$\pi_1(L(5,q))=\langle\mu\mid\mu^5=1\rangle\cong\Z/5$.
These are connected closed oriented manifolds. No homotopy-equivalence,
minimality, or all-rank classification assertion is needed.

\section{Exact certificates and limits}
The accompanying archive contains two full root-lattice checks using
Python's standard library and an independent full-weight modular-word
calculation using NumPy object arrays of arbitrary-precision integers.
The latter enumerates all 126 labels, 120 Weyl permutations and 8,001
symmetric matrix entries in $\Z[\zeta_{100}]$, reduced modulo
$\Phi_{100}=x^{40}-x^{30}+x^{20}-x^{10}+1$.
The lattice checks use $\Phi_{10}$ and $\Phi_{50}$ and enumerate all
625 quotient representatives; the fresh $\Phi_{50}$ route also checks
orientation, inverse-denominator, periodicity and $S^3$ boundaries.

Validation uses explicit exceptions that remain active under Python
\texttt{-O}. A manifest-bound runner records fresh child-process exit
codes, streams, source hashes and runtime information in both modes;
deliberately false arithmetic targets must fail in both modes. The
historical output field \texttt{exact\_assertions=2005} in one checker
counts its explicit finite checks; it is not a quality measure or 2,005
independent mathematical results. Repetition in two modes does not
double the distinct checks. Floating numbers in the fresh lattice
output are labeled diagnostics and decide no exact identity or sign.
The finite computations check this specialization; they do not prove
the general modular-category theorem or establish historical priority.

\section{Priority, assistance, and review}
Takahito Kuriya's \emph{The LMO invariant and the Guadagnini--Pilo
conjecture for lens spaces} is listed as a preprint by his Kyushu
institutional record \cite{K}. An institutional report records an
exact-title talk on February 21, 2003 \cite{KR}. Its full text and
complete conclusions and hypotheses could not be obtained. We cannot
establish that it does not already contain or circumscribe this result.
Historical priority therefore remains unresolved. Titles, aggregator
abstracts, later perturbative results, and bounded searches cannot
replace that comparison. No first counterexample, exhaustive novelty
clearance, or new historical resolution is claimed.

The separate priority supplement documents the read boundaries. In
particular, ordinary SU($N$) WZNW is explicitly spin-independent in the
later Kubo--Yokoyama comparison \cite{KY}; their proposed background
phases are a separate issue. Those $q=1$ tables do not establish an
earlier instance of this target pair. We make no blanket dismissal of
SU($N$) calculations as spin invariants.

AI tools were used extensively in solving, drafting, literature
comparison, exact reproduction, and adversarial verification. Three
separate AI validation families checked the arithmetic, normalization,
and root-lattice/topological argument. This note has not undergone
conventional human peer review and is unrefereed. The author takes
responsibility for the stated claims and qualifications. The original
author effort was two of five substantive approaches; preparation and
validation added no central proof-search turn.

\begin{thebibliography}{9}
\bibitem{O} T. Ohtsuki (ed.), \emph{Problems on invariants of knots and
3-manifolds}, Geometry \& Topology Monographs \textbf{4} (nominal 2002;
published June 1, 2004), 377--572; Conjecture 7.5, p.~474.
\url{https://msp.org/gtm/2002/04/gtm-2002-04-024s.pdf}.
\bibitem{GP} E. Guadagnini and L. Pilo, \emph{Three-manifold invariants
and their relation with the fundamental group}, Commun. Math. Phys.
\textbf{192} (1998), 47--65. DOI: \href{https://doi.org/10.1007/s002200050290}{10.1007/s002200050290}.
Available preprint: \href{https://arxiv.org/abs/hep-th/9612090v1}{arXiv:hep-th/9612090v1}.
\bibitem{HT} S. K. Hansen and T. Takata, \emph{Reshetikhin--Turaev
invariants of Seifert 3-manifolds for classical simple Lie algebras,
and their asymptotic expansions}, J. Knot Theory Ramifications
\textbf{13}(5) (2004), 617--668. DOI:
\href{https://doi.org/10.1142/S0218216504003342}{10.1142/S0218216504003342}.
Formula source used here: \href{https://arxiv.org/abs/math/0209403v2}{arXiv:math/0209403v2},
Eq.~(12), Sections 4--5, Theorem 5.1 (preprint p.~39).
\bibitem{K} T. Kuriya, \emph{The LMO invariant and the Guadagnini--Pilo
conjecture for lens spaces}, preprint, associated with a 2003 talk;
full text unavailable to this work. Institutional listing:
\url{https://www2.math.kyushu-u.ac.jp/coe/staff/staff_sub/kuriya.html}.
\bibitem{KR} Kyushu University 21st Century COE program, institutional
report, Kuriya activities, printed p.~121, recording the February 21,
2003 talk. \url{https://www2.math.kyushu-u.ac.jp/coe/report/pdf3/report-6.pdf}.
\bibitem{KY} N. Kubo and S. Yokoyama, \emph{Topological phase, spin
Chern--Simons theory and level rank duality on lens space}, JHEP
\textbf{04} (2022), 074. DOI:
\href{https://doi.org/10.1007/JHEP04(2022)074}{10.1007/JHEP04(2022)074};
\href{https://arxiv.org/abs/2108.09300v2}{arXiv:2108.09300v2}.
\end{thebibliography}
\end{document}
